% Generated draft: requires mathematical and source review.
\documentclass[11pt]{article}
\usepackage[margin=25mm]{geometry}
\usepackage{fontspec}
\setmainfont{Latin Modern Roman}
\usepackage{amsmath,amssymb,mathtools,unicode-math}
\setmathfont{Latin Modern Math}
\usepackage{enumitem,needspace,xurl,hyperref}
\setlength{\parindent}{0pt}
\setlength{\parskip}{5pt}
\setlength{\emergencystretch}{3em}
\newcommand{\N}{\mathbb{N}}
\newcommand{\Z}{\mathbb{Z}}
\newcommand{\Q}{\mathbb{Q}}
\newcommand{\R}{\mathbb{R}}
\newcommand{\C}{\mathbb{C}}
\newcommand{\F}{\mathbb{F}}
\newcommand{\A}{\mathbb{A}}
\newcommand{\PP}{\mathbb{P}}
\newcommand{\E}{\mathbb{E}}
\newcommand{\eps}{\varepsilon}
\newcommand{\dd}{\,\mathrm{d}}
\begin{document}
\section*{UnsolvedMath ID: 30000670}
\textbf{Source problem code:} catalogue entry
\subsection{Definitions and mathematical statement}
Retain the setting of the source's Theorem 2. Let \((K,D)\) be its Hasse--Schmidt field with absolute-constant field \(C\), and let \(A,B\) be commutative algebraic groups over \(C\). In a one-parameter Hasse--Schmidt convention, \(D=(D_i)_{i\geq0}\) has \(D_0=\operatorname{id}\) and
\[
 D_r(xy)=\sum_{i+j=r}D_i(x)D_j(y),\qquad
 C=\bigcap_{i\geq1}\ker D_i.
\]
The source does not specify iterativity or other field/derivation hypotheses. Let \(\phi:\widehat A\to\widehat B\) be the intended nowhere-algebraic formal isomorphism at the identities. The source writes \(\phi:A\to B\), but a formal map must not be silently treated as an algebraic morphism on all points.

For \(a\in A(K)\), \(b\in B(K)\), retain the differential equation and dimension condition
\[
 U_\phi(l_D^A(a))=l_D^B(b),\qquad
 \operatorname{trdeg}_C(a,b)\leq n.
\]
The logarithmic-derivative maps \(l_D^A,l_D^B\), induced map \(U_\phi\), nowhere-algebraic condition and integer \(n\) are not defined; in particular \(n\) is not explicitly identified with either group dimension. Under the complete hypotheses of the missing Theorem 2, the supplied record reports a conclusion
\[
 b\in B_0(K)+c_B,
 \qquad B_0<B\text{ a proper algebraic subgroup over }C,
 \quad c_B\in B(C).
\]
This is attributed to that theorem's omitted setup, not proved from the abbreviated displayed hypotheses alone. The three requests are as follows.
\begin{enumerate}
\item If \(\phi^{-1}\) is also nowhere algebraic, can the conclusion be strengthened on the \(A\)-side? The cleaned ``likewise'' reading asks for a proper algebraic subgroup \(A_0<A\) over \(C\) and a constant translate \(c_A\in A(C)\) with \(a\in A_0(K)+c_A\). The precise symmetry of the theorem's other hypotheses is not given, so application to the inverse is not asserted automatically.
\item Can the formal-map case from the multiplicative group \(\mathbb G_m\) to an ordinary elliptic curve prove weak \(\mathrm{CIT}_p\)? The original statement expressly says ``in characteristic zero,'' whereas the clean version drops that phrase and the title suggests positive characteristic. The meaning of ordinary, the relevant formal map and the exact weak \(\mathrm{CIT}_p\) assertion are absent. The characteristic-zero qualifier is retained as original-source data and the tension is flagged, not silently changed to characteristic \(p\).
\item What invariant should replace \(\dim_\Q\) in a Schanuel-style reformulation? The source gives neither the candidate vector space nor the proposed inequality, so no dimension over a different field or endomorphism ring is chosen.
\end{enumerate}
The literature triage checked 2026-08-21 reports no authoritative complete resolution found. Its open label and reported source verification remain dated metadata, not an independently verified theorem or a restoration of Theorem 2's missing assumptions.
\subsection{Short English statement}
The three tasks ask for a matching proper-subgroup coset on the first group, an application of the multiplicative-to-elliptic formal-map case to weak \(\mathrm{CIT}_p\), and the correct dimension invariant for a Schanuel formulation. The theorem setup and several key definitions are missing. The original characteristic-zero qualifier is preserved despite its omission from the cleaned question.
\subsection{Sources}
\begin{itemize}
\item Local UnsolvedMath registry, permanent ID 30000670, source code catalogue entry, ``Formal Group Isomorphisms and Positive-Characteristic CIT,'' including the three original Theorem 2 questions, the original characteristic-zero qualifier and expanded cleaned hypotheses.
\item \emph{Model Theory and Groups (2007)}, Oberwolfach Reports, pp.\ 83--138, DOI 10.4171/owr/2007/02: \url{https://doi.org/10.4171/owr/2007/02}. This is also the sole supplied literature reference.
\item Literature triage checked 2026-08-21; record updated 2026-09-16. No external source has been visited or independently checked here.
\end{itemize}
\end{document}
